\chapter{Game Context and Rules}
\label{annexe:two}

\fancyhead[LE]{\textsc{\chaptername~\thechapter}}

\section{Competitive Format}


Counter-Strike 2 (CS2) is a tactical, round-based first-person shooter played in 5-versus-5 matches between a Terrorist (T) and a Counter-Terrorist (CT) side. In the default competitive mode, a match consists of up to 24 rounds (12 per half before sides switch), each lasting a maximum of approximately 1 minute 55 seconds. A round is won either by completing an objective (planting and detonating the bomb for the T side, or defusing it / eliminating the opposing team for the CT side), by eliminating the entire opposing team, or by the timer running out.


Each player starts a match with a fixed amount of in-game currency and earns or loses money based on round outcomes, kills, and objective completions. This currency is spent during a "buy phase" at the start of each round on weapons, grenades (utility), and armor (Kevlar vest and helmet). This economic layer means that a player's loadout — and therefore their potential for damage output — varies from round to round, which is directly relevant to the telemetry features discussed in Chapter 3 (e.g. "premade lobby dependencies", utility efficiency).


\section{Core Mechanics Relevant to Cheating}

Several mechanical systems are central to understanding why the cheats described in Chapter 1 provide such a significant advantage:


\begin{itemize}
\item Recoil and spray patterns: each weapon follows a fixed, deterministic recoil pattern when fired continuously ("spraying"). Skilled players memorize and compensate for this pattern by moving their mouse in the opposite direction ("spray control"). An aimbot or recoil-compensation script can replicate this perfectly and instantly, which is the mechanism illustrated in Figure 2.
\item Movement and "counter-strafing": CS2 applies an accuracy penalty while a player is moving. Precise, frame-perfect counter-strafing (briefly tapping the opposite movement key to cancel momentum before shooting) restores full accuracy. This is a mechanical skill ceiling that automation cheats can trivially exceed.
\item Sound and information design ("fog of war"): as introduced in Section 1.1.2, large parts of player decision-making rely on incomplete information — footstep sounds, gunfire direction, utility usage, and map knowledge. The "mini-map" only shows teammates and recently-spotted enemies. This deliberate information asymmetry is the foundation that wallhacks and radarhacks subvert.
\item Damage model and hit zones: unlike some shooters, CS2 does not apply a fixed "headshot multiplier" as a flat bonus; rather, each of the four hit zones on a player model (head, chest/arms, stomach, legs) carries its own damage multiplier applied to the weapon's base damage, with the head dealing roughly four times the damage of a torso hit, the stomach somewhat more than the torso, and the legs somewhat less. This zone-dependent damage model is precisely why aim assistance (aimbots) is so impactful: it does not merely guarantee a hit, it guarantees a hit on the highest-value zone (the head), turning encounters that would normally require multiple shots into immediate eliminations.
\end{itemize}